% !TEX program = xelatex
\documentclass[11pt,a4paper]{article}

% ---- Layout and Typography ----
\usepackage{geometry}
\geometry{margin=2.5cm}
\usepackage{setspace}
\onehalfspacing
\usepackage{parskip}            % 段落间以空行分隔，无缩进

% ---- Fonts (XeLaTeX) ----
\usepackage{fontspec}
\usepackage{xeCJK}
% 按需设置你的中英文字体
% \setmainfont{Times New Roman}
% \setCJKmainfont{SimSun}

% ---- Micro-typography ----
\usepackage{microtype}

% ---- Math ----
\usepackage{amsmath,amssymb}

% ---- Tables ----
\usepackage{booktabs}
\usepackage{float}
\usepackage{array}          % 扩展表格列格式
\usepackage{caption}
\usepackage{subcaption}
\captionsetup{
    font=small,
    labelfont=bf,
    textfont=it,
    skip=6pt
}

% ---- Code Listings ----
\usepackage{listings}
\usepackage{xcolor}
\lstset{
    basicstyle=\ttfamily\small,
    columns=flexible,
    backgroundcolor=\color{gray!10},
    frame=single,
    breaklines=true,
    postbreak=\mbox{$\hookrightarrow$ },
    captionpos=b
}

% ---- Lists ----
\usepackage{enumitem}

% ---- Graphics ----
\usepackage{graphicx}

% ---- References ----
\usepackage{cite}

% ---- Hyperlinks ----
\usepackage{hyperref}
\hypersetup{
    colorlinks=true,
    linkcolor=blue,
    citecolor=blue,
    urlcolor=blue,
}

% ---- Title Area ----
\usepackage{titling}
\setlength{\droptitle}{-2em}   % 向上微调标题，避免开头太靠下

% ---- Title ----
\title{Classical Chinese Poetry Generation \\ with Transformer Architectures
    \\[0.5\baselineskip]
    \large From-Scratch Decoder-Only \& Encoder-Decoder vs.\
    Fine-Tuned GPT-2 Chinese Poem}
\author{He Yang \\[2pt] 
    \small\url{https://github.com/aSun-King-W/Poetry_generation}}
\date{May 10, 2026}

\begin{document}

\maketitle
\thispagestyle{empty}

% =====================================================================
\begin{abstract}
Classical Chinese poetry, with its strict tonal patterns and concise forms,
presents a unique challenge for natural language generation.
This project explores three Transformer-based architectures for the
couplet-completion task: given an upper verse line (上句), generate a
semantically and metrically coherent lower line (下句).
We implement a Decoder-Only Transformer and an Encoder-Decoder Transformer
entirely from scratch using basic PyTorch tensor operations, and compare them
against a fine-tuned pretrained model (\texttt{uer/gpt2-chinese-poem}).
All models are trained on over one million couplets extracted from the complete
Tang and Song poetry corpora.
The pretrained model achieves the lowest perplexity (1.49), while the
from-scratch models demonstrate competitive learning with 11M--12M parameters
(PPL 79.15 and 102.44) and provide valuable insights into architectural
trade-offs.
We further deploy a Gradio-based interactive demo for qualitative side-by-side
comparison.
\end{abstract}

% =====================================================================
\section{Introduction}
\label{sec:intro}

Classical Chinese poetry (古典诗歌) is a cultural heritage spanning thousands of
years. Forms such as \emph{five-character quatrain} (五言绝句) and
\emph{seven-character regulated verse} (七言律诗) impose strict constraints on
character count, tonal alternation, and parallel couplet structure.
The couplet-completion task --- given the first line of a couplet, generate the
second --- tests a model's ability to capture both semantic coherence and formal
poetic constraints.

Neural approaches to poetry generation have evolved from RNN-based
methods~\cite{zhang2017chinese, yi2017generating} to Transformer-based
architectures~\cite{vaswani2017attention}.
The release of the \texttt{uer/gpt2-chinese-poem} model~\cite{zhao2019uer},
pretrained on 800K classical Chinese poems, provides a strong baseline for
fine-tuning.
However, from-scratch implementations remain valuable for understanding
architectural decisions and for deployment in resource-constrained environments.

In this project, we implement three distinct approaches:
\begin{enumerate}[nosep]
    \item \textbf{Decoder-Only Transformer:} A causal language model trained
          to predict the full sequence ``upper verse + [SEP] + lower verse''
          autoregressively.
    \item \textbf{Encoder-Decoder Transformer:} A sequence-to-sequence model
          where an encoder processes the upper verse bidirectionally and a
          decoder generates the lower verse with cross-attention.
    \item \textbf{Pretrained GPT-2 Fine-Tuning:}
          The \texttt{uer/gpt2-chinese-poem} model, adapted to the
          couplet-completion task via continued training.
\end{enumerate}

All Transformer components for methods (1) and (2) are implemented from
scratch using only basic PyTorch tensor operations, without relying on
built-in Transformer APIs such as \texttt{torch.nn.Transformer} or
\texttt{torch.nn.MultiheadAttention}.

% =====================================================================
\section{Methodology}
\label{sec:method}

\subsection{Dataset and Preprocessing}

We use the \texttt{chinese-poetry/chinese-poetry}
corpus~\cite{chinese-poetry}, which contains the complete
\emph{Quan Tang Shi} (全唐诗, Complete Tang Poems) and
\emph{Song Shi} (宋诗, Song Poems).
The preprocessing pipeline consists of:

\begin{enumerate}[nosep]
    \item Extract individual lines from each poem's \texttt{paragraphs} field.
    \item Filter to lines of exactly 5 or 7 Chinese characters (the standard
          verse lengths).
    \item Pair adjacent lines within the same poem: lines 0--1 form the first
          couplet, lines 2--3 the second, etc.
          (quatrains yield 2 pairs; regulated verses yield 4).
    \item Retain only pairs where both lines have the same character length.
\end{enumerate}

The final dataset comprises \textbf{1,037,826} five-character couplet pairs.
The data is randomly partitioned into training (80\%), validation (10\%), and
test (10\%) sets.

A character-level vocabulary is constructed by counting character frequencies
across all pairs.
Characters appearing fewer than 3 times are mapped to \texttt{<UNK>}.
The resulting vocabulary contains \textbf{9,119} tokens, including 5 special
tokens: \texttt{<PAD>}, \texttt{<UNK>}, \texttt{<BOS>}, \texttt{<EOS>},
and \texttt{<SEP>}.

\subsection{Data Formats for Each Architecture}

Each model architecture requires a different input--target format:

\begin{itemize}[nosep]
    \item \textbf{Decoder-Only:}
          The input sequence is ``upper verse + SEP + lower verse + EOS''.
          During training, the model predicts the entire sequence shifted right
          by one position under a causal attention mask.
          The SEP token signals the boundary between input and target.
    \item \textbf{Encoder-Decoder:}
          The encoder receives the upper verse alone.
          The decoder receives ``BOS + lower verse'' as input and is trained
          to predict ``lower verse + EOS''.
          Cross-attention in each decoder block attends to the encoder's
          output.
    \item \textbf{Pretrained:}
          The \texttt{BertTokenizer} (required by
          \texttt{uer/gpt2-chinese-poem}) expects spaces between Chinese
          characters.
          Inputs are formatted as ``[CLS] u p p e r [SEP] l o w e r''.
          Labels mask out the prefix (up to and including [SEP]) so the model
          only learns to predict the lower verse.
\end{itemize}

\subsection{From-Scratch Transformer Components}
\label{sec:components}

All core Transformer components are implemented in
\texttt{models/transformer.py} using only \texttt{torch.Tensor} operations,
\texttt{nn.Linear}, \texttt{nn.Embedding}, \texttt{nn.Dropout}, and
\texttt{F.gelu}.
The key components are:

\subsubsection*{Scaled Dot-Product Multi-Head Attention}
Given queries \(Q\), keys \(K\), and values \(V\):
\[
\text{Attention}(Q,K,V) =
    \text{softmax}\!\left(\frac{QK^{\mathsf{T}}}{\sqrt{d_k}}\right)V
\]
We project \(Q\), \(K\), \(V\) through learned linear layers, split into
\(n_{\text{heads}}=8\) heads each of dimension
\(d_k = d_{\text{model}} / n_{\text{heads}}\), compute attention scores
independently per head, concatenate outputs, and apply a final projection
\(W_O\).
Causal masking (for decoder self-attention) is implemented as a
lower-triangular binary mask; padding masks are combined via element-wise
multiplication.

\subsubsection*{Sinusoidal Positional Encoding}
Position information is injected via fixed sinusoidal encodings:
\[
\begin{aligned}
PE(pos, 2i)   &= \sin\!\left(pos / 10000^{2i/d_{\text{model}}}\right) \\
PE(pos, 2i+1) &= \cos\!\left(pos / 10000^{2i/d_{\text{model}}}\right)
\end{aligned}
\]
The encoding supports sequences up to length 256 and is not learned.

\subsubsection*{Feed-Forward Network}
Each position is processed through a two-layer network:
\(d_{\text{model}} \to d_{\text{ff}} \to d_{\text{model}}\) with GELU
activation and dropout.

\subsubsection*{Pre-Norm Residual Architecture}
Each sublayer (self-attention, cross-attention, feed-forward) follows a
Pre-Norm design: LayerNorm is applied \emph{before} the sublayer, followed
by a residual connection and dropout.
This has been shown to stabilize training in deep
transformers~\cite{xiong2020layer}.

\subsubsection*{Decoder and Encoder Blocks}
The \texttt{TransformerBlock} (decoder) contains three sublayers:
masked self-attention, cross-attention, and feed-forward.
The \texttt{EncoderBlock} contains two sublayers:
bidirectional self-attention and feed-forward.

\subsection{Architecture Details}
\label{sec:architectures}

\subsubsection*{Decoder-Only Transformer (\texttt{models/decoder\_only.py})}
\[
\text{Embedding} \to \text{PE} \to
    [\text{TransformerBlock}]^{\times N} \to
    \text{LayerNorm} \to \text{Linear}(d_{\text{model}}, V)
\]
where \(V = 9119\) is the vocabulary size.
Configuration: \(d_{\text{model}}=256\), \(N=6\), \(n_{\text{heads}}=8\),
\(d_{\text{ff}}=1024\), dropout=0.1, \(\text{max\_len}=32\).
Total parameters: \textbf{11.0M}.

\subsubsection*{Encoder-Decoder Transformer (\texttt{models/encoder\_decoder.py})}
The encoder uses \(N_{\text{enc}}=4\) \texttt{EncoderBlock}s;
the decoder uses \(N_{\text{dec}}=4\) \texttt{TransformerBlock}s with
cross-attention attending to encoder outputs.
Embeddings and positional encodings are shared between encoder and decoder.
Configuration: \(d_{\text{model}}=256\), \(N_{\text{enc}}=4\),
\(N_{\text{dec}}=4\), \(n_{\text{heads}}=8\), \(d_{\text{ff}}=1024\),
dropout=0.1, \(\text{max\_len}=32\).
Total parameters: \textbf{12.1M}.

\subsubsection*{Pretrained GPT-2 Chinese Poem (\texttt{models/pretrained.py})}
We load \texttt{uer/gpt2-chinese-poem}, a GPT-2 model with
\(d_{\text{model}}=768\), 12 layers, 12 attention heads, and
\(d_{\text{ff}}=3072\) (approximately 102M parameters).
The model was originally pretrained on 800K classical Chinese poems using
masked language modeling.
We fine-tune all parameters (full fine-tuning, not adapter-based) on our
couplet dataset.

\subsection{Training Setup}

All three models are trained with the Adam optimizer.
The from-scratch models use \(\beta_1=0.9\), \(\beta_2=0.999\),
\(\epsilon=10^{-8}\), a learning rate of \(10^{-4}\) with 2,000 warmup steps
followed by cosine decay to zero, gradient clipping at \(\text{norm}=1.0\),
and label smoothing \(\epsilon=0.1\).
The pretrained model uses AdamW with weight decay 0.01, a lower learning rate
of \(2\times 10^{-5}\), and 500 warmup steps.
Cross-entropy loss ignores padding positions (set to \(-100\)).

\begin{table}[H]
\centering
\small
\renewcommand{\arraystretch}{1.1}
\caption{Hyperparameter comparison across the three models.}
\label{tab:hyperparams}
\begin{tabular}{lcccc}
\toprule
\textbf{Hyperparameter} & \textbf{Dec-Only} & \textbf{Enc-Dec} & \textbf{Pretrained} \\
\midrule
\(d_{\text{model}}\)    & 256   & 256   & 768 \\
Layers                  & 6     & 4 + 4 & 12 \\
Attention heads         & 8     & 8     & 12 \\
\(d_{\text{ff}}\)       & 1024  & 1024  & 3072 \\
Dropout                 & 0.1   & 0.1   & 0.1 \\
Learning rate           & \(10^{-4}\) & \(10^{-4}\) & \(2\times10^{-5}\) \\
Batch size              & 256   & 1024  & 64 \\
Max sequence length     & 32    & 32    & 64 \\
Epochs                  & 30    & 30    & 10 \\
Weight decay            & 0     & 0     & 0.01 \\
Parameters              & 11.0M & 12.1M & \(\sim\)102M \\
\bottomrule
\end{tabular}
\end{table}

\subsection{Decoding Strategies}

During inference, we support three decoding strategies:
\begin{enumerate}[nosep]
    \item \textbf{Greedy decoding:} At each step, select the token with the
          highest probability.
    \item \textbf{Beam search:} Maintain \(k=5\) partial hypotheses, selecting
          the top-\(k\) sequences by cumulative log-probability at each step.
    \item \textbf{Temperature sampling:} Sample from the token distribution
          scaled by temperature \(t \in [0.1, 2.0]\), with \(t=0.8\) as
          default.
\end{enumerate}

% =====================================================================
\section{Experiments}
\label{sec:experiments}

\subsection{Evaluation Metrics}

We evaluate models using two automatic metrics:
\begin{itemize}[nosep]
    \item \textbf{Perplexity (PPL):}
          \(\exp(\mathcal{L}_{\text{test}})\), where
          \(\mathcal{L}_{\text{test}}\) is the average cross-entropy loss on
          the test set.
          Lower perplexity indicates better predictive performance.
    \item \textbf{BLEU-4:}
          The 4-gram overlap between the generated lower verse and the
          reference lower verse, computed using the corpus BLEU
          method~\cite{papineni2002bleu} with additive smoothing.
\end{itemize}

\subsection{Results}

\begin{table}[H]
\centering
\small
\renewcommand{\arraystretch}{1.1}
\caption{Evaluation results on the test set.
         Perplexity is computed on the full test set; BLEU-4 is computed
         over a subset of 500 samples using greedy decoding.}
\label{tab:results}
\begin{tabular}{lcc}
\toprule
\textbf{Model}                          & \textbf{Perplexity \(\downarrow\)} & \textbf{BLEU-4 \(\uparrow\)} \\
\midrule
Decoder-Only (from scratch)             & 79.15   & 0.0010 \\
Encoder-Decoder (from scratch)          & 102.44  & 0.0010 \\
Pretrained GPT-2 (fine-tuned)           & 1.49    & 0.0136 \\
\bottomrule
\end{tabular}
\end{table}

Table~\ref{tab:results} presents the main quantitative results.
The pretrained GPT-2 model dramatically outperforms both from-scratch models
in perplexity (1.49 vs.~79.15 and 102.44), which is expected given its much
larger capacity (102M vs.~11M parameters) and substantial pretraining on
related poetry data.
The BLEU-4 scores follow the same ordering, though absolute values are low
across all models --- a known phenomenon in creative text generation where
many valid continuations exist for a given
prompt~\cite{novikova2017why}.

Among the two from-scratch models, the Decoder-Only architecture achieves
better perplexity (79.15) than the Encoder-Decoder (102.44).
This is somewhat surprising given that the Encoder-Decoder has slightly more
parameters and bidirectional encoding.
We attribute this to two factors:
(1) the Decoder-Only model sees the upper verse as part of its autoregressive
context, while the Encoder-Decoder must learn to attend through
cross-attention;
(2) the default hyperparameters (layer counts) may not be optimally balanced
for the Encoder-Decoder.

\subsection{Qualitative Analysis}

\begin{table}[H]
\centering
\small
\renewcommand{\arraystretch}{1.2}
\caption{Generation examples from all three models using beam search
         (beam=5).}
\label{tab:examples}
\begin{tabular}{p{2.5cm}p{3.2cm}p{3.2cm}p{3.2cm}}
\toprule
\textbf{Input (上句)} & \textbf{Decoder-Only} & \textbf{Encoder-Decoder} & \textbf{Pretrained GPT-2} \\
\midrule
床前明月光      & 枕上清風起    & 窗外清風清    & 下寒泉聲 \\
举头望明月      & 不知何處來    & 飛鴻飛上天    & 低頭拾落花 \\
白日依山尽      & 青山带水流    & 青雲入水涯    & 黄花逐水流 \\
春眠不觉晓      & 夜夢不成眠    & 夜夢不可聽    & 春夢有誰知 \\
海内存知己      & 江湖失故人    & 山川不見人    & 天涯失友生 \\
\bottomrule
\end{tabular}
\end{table}

Table~\ref{tab:examples} shows qualitative generation results for classic
upper verses.
With limited capacity (\textasciitilde{}11M--12M parameters), the from-scratch
models do not reproduce the original ground-truth lower verses, but generate
semantically reasonable and contextually coherent alternatives.
The pretrained GPT-2 produces more metrically regular and antithetical output,
demonstrating the advantage of large-scale pretraining.

\subsection{Training Dynamics}

The from-scratch models exhibit gradual convergence over 30 epochs.
The pretrained model converged much faster (10 epochs) thanks to its strong
language model prior.

\subsection{Implementation Challenges}

Several challenges arose during implementation:

\begin{enumerate}[nosep]
    \item \textbf{Padding mask combination:}
          Combining the causal mask and padding mask in the Decoder-Only model
          required careful broadcasting (\texttt{causal\_mask * pad\_mask})
          to ensure that padding positions receive \(-\infty\) attention
          scores.
    \item \textbf{Cross-attention in the Encoder-Decoder:}
          The encoder output must be broadcast to each beam during beam search
          decoding, requiring explicit dimension expansion with
          \texttt{.repeat()}.
    \item \textbf{Pretrained tokenizer spacing:}
          The \texttt{BertTokenizer} for \texttt{uer/gpt2-chinese-poem}
          requires spaces between Chinese characters, which is non-standard
          for Chinese text processing.
    \item \textbf{DataParallel checkpoint compatibility:}
          Models trained with multi-GPU DataParallel save state dicts with a
          \texttt{module.} prefix; loading on a single GPU requires key
          remapping.
\end{enumerate}

% =====================================================================
\section{Conclusion}
\label{sec:conclusion}

We presented a systematic comparison of three Transformer architectures for
classical Chinese poetry couplet completion.
Our main findings are:

\begin{enumerate}[nosep]
    \item From-scratch Transformer implementations with modest capacity
          (11M--12M parameters) can learn meaningful poetry generation
          capabilities, achieving reasonable perplexity and generating
          semantically coherent lower verses.
    \item The Decoder-Only architecture outperforms the Encoder-Decoder in
          our setting (PPL 79.15 vs.~102.44), likely due to more effective
          use of the upper-verse context and better-tuned capacity allocation.
    \item Fine-tuning a domain-specific pretrained model
          (uer/gpt2-chinese-poem) yields dramatically superior results
          (PPL 1.49 vs.~79--102), demonstrating the value of large-scale
          pretraining even for specialized literary domains.
\end{enumerate}

\textbf{Limitations.}
The from-scratch models, while instructive, are far from matching the quality
of pretrained approaches.
Their relatively high perplexity indicates that
(1) 11M parameters may be insufficient for capturing the full complexity of
Chinese poetic form, and
(2) architectural improvements such as learned positional encodings or deeper
layers could bridge part of the gap.

\textbf{Future work.}
Several directions are promising:
(a) scaling from-scratch models to match pretrained capacity;
(b) incorporating explicit poetic constraints (tone patterns, rhyme) into the
loss function;
(c) exploring adapter-based fine-tuning for parameter-efficient adaptation;
(d) applying reinforcement learning from human feedback to improve generation
quality;
(e) extending the task to full poem generation beyond single couplets.

\begin{thebibliography}{10}

\bibitem{vaswani2017attention}
A.~Vaswani, N.~Shazeer, N.~Parmar, \emph{et al.},
``Attention is all you need,''
in \emph{Advances in Neural Information Processing Systems}, 2017.

\bibitem{zhang2017chinese}
X.~Zhang and M.~Lapata,
``Chinese poetry generation with recurrent neural networks,''
in \emph{Proceedings of EMNLP}, 2017.

\bibitem{yi2017generating}
X.~Yi, R.~Li, and M.~Sun,
``Generating Chinese classical poems with RNN encoder-decoder,''
in \emph{Chinese Computational Linguistics and Natural Language Processing}, 2017.

\bibitem{zhao2019uer}
Z.~Zhao, H.~Chen, J.~Zhang, \emph{et al.},
``UER: An open-source toolkit for pre-training models,''
in \emph{Proceedings of EMNLP-IJCNLP: System Demonstrations}, 2019.

\bibitem{xiong2020layer}
R.~Xiong, Y.~Yang, D.~He, \emph{et al.},
``On layer normalization in the transformer architecture,''
in \emph{Proceedings of ICML}, 2020.

\bibitem{papineni2002bleu}
K.~Papineni, S.~Roukos, T.~Ward, and W.-J.~Zhu,
``BLEU: a method for automatic evaluation of machine translation,''
in \emph{Proceedings of ACL}, 2002.

\bibitem{novikova2017why}
J.~Novikova, O.~Dušek, A.~Cercas Curry, and V.~Rieser,
``Why we need new evaluation metrics for NLG,''
in \emph{Proceedings of EMNLP}, 2017.

\bibitem{chinese-poetry}
\texttt{chinese-poetry/chinese-poetry},
``The most comprehensive database of Chinese poetry,''
\url{https://github.com/chinese-poetry/chinese-poetry}.

\end{thebibliography}

\end{document}